\chapter{Projekt systemu}
\label{cha:projekt_systemu}

Projekt systemu zaprojektowano z~myślą o~zadaniach konkursowych EvSLAM i~M3ED. Zbiory danych różnią się sposobem organizacji informacji, strukturą pomiarów oraz zakresem dostępnej kalibracji, dlatego wymagają odrębnych etapów wczytania i~wstępnego przygotowania. Po ujednoliceniu danych wejściowych dalsze przetwarzanie realizuje wspólny algorytm estymacji ruchu i~budowy mapy.

W~rozdziale przedstawiono wymagania wynikające z~obu konkursów, przyjęty zestaw czujników oraz założenia projektowe systemu. Następnie opisano strukturę rozwiązania, przepływ danych pomiędzy modułami oraz najważniejsze elementy odpowiedzialne za estymację ruchu, mapowanie i~korektę trajektorii.


\section{Założenia wynikające z konkursów}
\label{sec:zalozenia_systemu}

\subsection{EvSLAM Challenge}
\label{subsec:projekt_evslam}

EvSLAM Challenge dotyczy estymacji stanu podczas szybkich manewrów. Udostępnione dane obejmują dwa strumienie z kamer zdarzeniowych DVXplorer, pomiary IMU oraz obrazy RGB, monochromatyczne obrazy z kamer podczerwieni i mapy głębi rejestrowane przez kamerę Intel RealSense D435. Kamery zdarzeniowe mają rozdzielczość $640\times480$ i~bazę $10\mathrm{cm}$, natomiast IMU pracuje z~częstotliwością $800\mathrm{Hz}$~\cite{evslam2025data}. Rozmieszczenie czujników na platformie przedstawiono na rysunku~\ref{fig:evslam_sensors}. Konkurs przewiduje wariant wykorzystujący wyłącznie kamery zdarzeniowe jako źródło informacji wizyjnej oraz wariant łączący zdarzenia z~monochromatycznymi obrazami z kamery Intel RealSense D435~\cite{evslam2025challenge}.

\begin{figure}[!ht]
\centering
\includegraphics[width=0.9\linewidth]{evslam_sensors.png}
\caption{Zestaw czujników wykorzystywany do rejestracji danych EvSLAM na platformie czterowirnikowego drona. Źródło: \cite{evslam2025data}.}
\label{fig:evslam_sensors}
\end{figure}

Dla każdej sekwencji należy wyznaczyć w~zadanych chwilach pozycję i~orientację lewej kamery oraz jej prędkość liniową wyrażoną w~układzie kamery. Ocena uwzględnia bezwzględny błąd trajektorii (ang. \textit{Absolute Trajectory Error}, ATE) oraz pole pod krzywą współczynnika sukcesu dla kolejnych progów względnego błędu prędkości (ang. \textit{Area Under the Curve}, AUC). Pierwsza z tych metryk jest brana pod uwagę wyłącznie w przypadku, w którym druga (AUC), jest równa dal dwóch różnych algorytmów.

\subsection{M3ED SLAM Challenge}
\label{subsec: project_m3ed}

Zbiór M3ED został zarejestrowany z~wykorzystaniem rozbudowanego zestawu czujników obejmującego stereoskopową parę kamer zdarzeniowych Prophesee EVK4 o~rozdzielczości $1280\times720$, stereoskopową parę kamer rejestrujących obrazy w~skali szarości, kamerę RGB, jednostkę IMU oraz 64-kanałowy skaner laserowy~\cite{m3ed2025data}. Dane pochodzą z~różnych platform, obejmujących samochód, bezzałogowy statek powietrzny oraz robota kroczącego, dzięki czemu zawierają ruch o~znacznie różniącej się dynamice i~charakterze.

W~M3ED SLAM Challenge wykorzystanie czujników jest ograniczone do dwóch wariantów. W~wariancie \textit{Event (+ IMU)} pozycja może być wyznaczana na podstawie jednej lub dwóch kamer zdarzeniowych, opcjonalnie z~wykorzystaniem IMU. Wariant \textit{Event + Mono (+ IMU)} dopuszcza dodatkowo obrazy z~jednej kamery globalnej migawki (ang. \textit{global shutter camera})~\cite{m3ed2025challenge}.

Konkurs obejmuje trzy sekwencje testowe: jazdę samochodem w~środowisku miejskim, szybki lot UAV oraz ruch robota kroczącego. Dla wskazanych znaczników czasu należy wyznaczyć pozycję i~orientację referencyjnej kamery zdarzeniowej. Dokładność jest oceniana za pomocą bezwzględnego błędu pozy (ang. \textit{Absolute Pose Error}, APE) względem trajektorii referencyjnej~\cite{m3ed2025challenge}. W~przeciwieństwie do EvSLAM wynik nie obejmuje prędkości liniowej -- wymaganym rezultatem jest wyłącznie trajektoria położenia robota.


\subsection{Wybór konfiguracji czujników}
\label{subsec:wybor_czujnikow}

W~pracy przyjęto wspólny dla obu konkursów wariant wykorzystujący
\textbf{dwie kamery zdarzeniowe i~IMU}. Informacja wizyjna pochodzi wyłącznie
ze zdarzeń, dzięki czemu system odpowiada wariantom niewymagającym obrazów z kamer klatkowych. Konfiguracja stereo została wybrana z~dwóch powodów:

\begin{itemize}
    \item znana baza kamer pozwala od razu wyznaczać głębokość i~metryczną skalę translacji;
    \item dwie obserwacje tej samej chwili ograniczają niejednoznaczność geometrii charakterystyczną dla pojedynczej kamery.
\end{itemize}

Z~IMU wykorzystywana jest prędkość kątowa. Żyroskop pomaga ograniczyć wpływ
szybkich obrotów na reprezentację zdarzeń oraz może dostarczyć pomocniczą
informację o~zmianie orientacji, podobnie jak w~systemach łączących zdarzenia
z~pomiarami inercyjnymi~\cite{rebecq2017eventvio,vidal2018ultimate}.
Akcelerometr nie jest uwzględniany, dlatego opracowane rozwiązanie nie jest
pełnym systemem VIO: nie estymuje obciążeń pomiarowych IMU
ani wspólnego stanu wizyjno-inercyjnego. Jest to świadome ograniczenie zakresu,
pozwalające skupić projekt na geometrii zdarzeniowego stereo VO i~SLAM.

\section{Cele i organizacja systemu}
\label{sec:architektura_systemu}

System zaprojektowano tak, aby po przygotowaniu danych wejściowych zbiory EvSLAM i~M3ED mogły być przetwarzane przez wspólną część algorytmu. Najważniejsze cele projektowe obejmują:

\begin{itemize}
\item estymację trajektorii referencyjnej kamery zdarzeniowej w~skali metrycznej;
\item wyznaczenie prędkości liniowej z~trajektorii na potrzeby zadania EvSLAM;
\item możliwość pracy wyłącznie jako odometria wizyjna lub z~dodatkowym wykorzystaniem mapy;
\item relokalizację po utracie śledzenia oraz korektę dryfu po wykryciu ponownie odwiedzonego miejsca.
\end{itemize}

Ogólną organizację systemu i~przepływ danych przedstawiono na rysunku~\ref{fig:system_architecture}. Główny tor przetwarzania prowadzi od zsynchronizowanych strumieni zdarzeń stereo do trajektorii lewej kamery i~wyznaczanej z~niej prędkości liniowej. Opcjonalne pomiary żyroskopu mogą wspierać przygotowanie obserwacji oraz estymację pozy, natomiast wariant SLAM rozszerza system o~rzadką mapę, rozpoznawanie miejsc, relokalizację i~korektę grafu pozy.

\begin{figure}[!ht]
    \centering
    \includegraphics[width=1\linewidth]{system_architecture.pdf}
    \caption{Architektura i~przepływ danych w~opracowanym systemie. Zielone pola oznaczają dane wejściowe, pomarańczowe wyniki, a~niebieskie elementy wariantu SLAM. Przerywana strzałka przedstawia warunkowe odzyskanie pozy przez relokalizację.}
    \label{fig:system_architecture}
\end{figure}

Przepływ danych w~systemie przedstawiono na rysunku~\ref{fig:system_architecture}. Zdarzenia z~obu kamer są synchronizowane i~dzielone na wspólne okna czasowe, a~następnie poddawane filtracji oraz korekcji zniekształceń wywołanych obrotem kamery. Wykorzystywane są przy tym zsynchronizowane z~kamerami pomiary żyroskopu, natomiast pomiary akcelerometru nie uczestniczą w~obliczeniach. Przygotowane zdarzenia są przekształcane w~reprezentacje obrazowe i~rektyfikowane. System wykrywa na nich punkty, śledzi je w~czasie i~wyznacza ich głębokość na podstawie obserwacji z~obu kamer. Ruch kamery jest określany z~wykorzystaniem punktów 3D z~poprzedniego okna lub mapy oraz ich bieżących obserwacji; w~konfiguracji końcowej translacja pochodzi z~wizji, a~zmiana orientacji z~żyroskopu. Wybrane obserwacje rozwijają rzadką mapę, która wspiera dalszą lokalizację, rozpoznawanie wcześniej odwiedzonych miejsc i~korektę trajektorii. Wynikiem jest trajektoria lewej kamery, a~dla EvSLAM również wyznaczona na jej podstawie prędkość liniowa. Dane referencyjne nie uczestniczą w~estymacji i~są wykorzystywane wyłącznie podczas oceny wyników.
